

\vspace*{0.35\textheight}
\begin{flushright}
\itshape
À Henri Burnage,\\
mon Papili.
\end{flushright}

\newpage 
\chapter*{Acknowledgements}
    \par
    La premier contact LKB je crois que c'était fin 2019-début 2020. Donc ça commence à faire un petit bout de temps quand même, j'ai rencontré plein de gens, j'ai appris plein de choses, j'ai vécu plein de choses, et j'ai passé plein de bons moments, de moins bons aussi. Je tiens à remercier tous ceux qui ont contribué à faire de cette expérience une expérience aussi enrichissante, aussi bien sur le plan scientifique que sur le plan humain. Je vais essayer de faire une liste aussi exhaustive que possible, mais je vais oublier plein de gens, et je vais pas écrire un chapitre entier non plus (il est 2h du mat', le jour de la deadline). 

    \noindent \textbf{le LKB}
    Merci Pierre-François de m'avoir accueilli au LKB, et de m'avoir donné la chance de faire ma thèse dans un environnement aussi stimulant et bienveillant. Tu m'as toujours donné les clé
    
    
    Merci Antoine, pour  
    et nous fournir les moyens de faire de la recherche de pointe dans un environnement aussi stimulant et bienveillant. 
    \noindent \textbf{H24} \\
    \noindent \textbf{les copains LKB} \\
    \noindent \textbf{les copains d'ailleurs} \\
    \noindent \textbf{les innomables} \\
    \noindent \textbf{la famille} \\


    Ma venue au LKB ne tient pas à grand chose. En Décembre 2019, je suis le cours d'optique quantique de Michael Vanner au chaud dans un amphi londonien quand il mentionne l'optomécanique quantique. Dans la foulée je me renseigne sur le sujet, lis un peu de biblio et envoie un mail à Antoine pour lui demander s'il recrute des stagiaires. 

    Je tiens à remercier Antoine Heidmann et Nicolas Treps, les deux directeurs du labo. mes deux N+2 d'abord en thèse, puis en postdoc, 

